\documentclass{tufte-handout}

%\geometry{showframe}% for debugging purposes -- displays the margins

\usepackage{amsmath}
\usepackage{amssymb}
\usepackage{amsthm}
\usepackage{thmtools}
\declaretheorem[numberwithin=section]
{theorem, definition}
\declaretheorem{lemma, proposition, corollary}[
style=plain,
numberwithin=theorem
]
% Set up the images/graphics package
\usepackage{graphicx}
\setkeys{Gin}{width=\linewidth,totalheight=\textheight,keepaspectratio}
\graphicspath{{graphics/}}

\title{Lecture 3 MATH 576}
\author{Shan Gao}
\date{Fall 2023}  % if the \date{} command is left out, the current date will be used

% The following package makes prettier tables.  We're all about the bling!
\usepackage{booktabs}

% The units package provides nice, non-stacked fractions and better spacing
% for units.
\usepackage{units}

% The fancyvrb package lets us customize the formatting of verbatim
% environments.  We use a slightly smaller font.
\usepackage{fancyvrb}
\fvset{fontsize=\normalsize}

% Small sections of multiple columns
\usepackage{multicol}

% Provides paragraphs of dummy text
\usepackage{lipsum}

% These commands are used to pretty-print LaTeX commands
\newcommand{\doccmd}[1]{\texttt{\textbackslash#1}}% command name -- adds backslash automatically
\newcommand{\docopt}[1]{\ensuremath{\langle}\textrm{\textit{#1}}\ensuremath{\rangle}}% optional command argument
\newcommand{\docarg}[1]{\textrm{\textit{#1}}}% (required) command argument
\newenvironment{docspec}{\begin{quote}\noindent}{\end{quote}}% command specification environment
\newcommand{\docenv}[1]{\textsf{#1}}% environment name
\newcommand{\docpkg}[1]{\texttt{#1}}% package name
\newcommand{\doccls}[1]{\texttt{#1}}% document class name
\newcommand{\docclsopt}[1]{\texttt{#1}}% document class option name

\begin{document}
\section{Continuation of Hausdorff and etc}
\begin{proposition}[Limits of sequeneces in Hausdorff space are unique]
    IF $X$ is Hausdorff, then a sequence in $X$ has at most one limit
    
\end{proposition}
\begin{proof}
    Suppose $\{x_n\}_{n \in \mathbb{N}}$ converges to $x$ and $y$ with $x \neq y$. By $T_2$, we have that there exists $U$ and $V$ open in $X$ such that $x \in U$ and $y \in V$ and $X \cap V = \emptyset$. Then that means that there exists $N_1$ such that if $n \geq N_1$, then $x_n \in U$, similarly $\exists N_2$ such that if $n \geq N_2$, $x_n \in V$. Then pick $N = \max(N_1, N_2)$, we have that $x_n \in U$ and $x_n \in V$, which contradicts the fact that $U \cap V = \emptyset$. 
\end{proof}
\begin{definition}[basis]
    let $X$ be a topological space and $\mathcal{B}$ is called a basis if it's a collection of sets such that for all $U \in X$ open in $X$, $U$ is a union of elements in $\mathcal{B}$, so 
    \[
        U = \bigcup_{B \subseteq U \colon B \in \mathcal{B}}
        \]
\end{definition}


\end{document}